\documentclass[11pt]{article}

% --- Packages ---
\usepackage[margin=1in]{geometry}
\usepackage[hidelinks]{hyperref}
\usepackage{setspace}
\usepackage{parskip}

% --- Formatting ---
\setlength{\parskip}{0.5em}
\setlength{\parindent}{0em}

\begin{document}

\begin{center}
{\Large \textbf{Closed-Loop Bayesian Optimization of Multi-Material\\
3D-Printed Tensegrity-Inspired Energy Absorbers}}

\vspace{1em}
\textit{ASME IDETC-CIE 2026 --- Presentation-Only Abstract}

\vspace{0.5em}
\textit{Track: DAC-10 --- Design of Engineering Materials and Structures}
\end{center}

\vspace{1em}

Tensegrity-inspired architectures, in which rigid struts are suspended
within a continuous flexible network, can exhibit tunable nonlinear
mechanical responses and high strength-to-weight ratios, motivating
their use in lightweight protective and energy-absorbing structures.
Multi-material fused deposition modeling (FDM) can co-print rigid (PLA)
struts and flexible (TPU) elements in a single build, but the resulting
design space---spanning strut geometry, tension-element cross-section,
connectivity topology, and unit-cell tiling---is too large to explore
efficiently by trial-and-error experimentation.

This work presents a closed-loop experimental campaign that uses
Bayesian optimization (BO) to develop multi-material 3D-printed
tensegrity-inspired structures for improved energy absorption under
quasi-static compression and drop-weight impact loading.  The key
contribution is a design--print--test workflow in which experimental
measurements are used directly to update a Gaussian-process surrogate
model, enabling sequential design selection without relying on
finite-element simulation.  The surrogate is trained on peak transmitted
force, specific energy absorption (SEA), and compaction efficiency, and
a batch acquisition function recommends the next set of candidate
designs after each test round.  The workflow combines automated model
updating and experiment selection with manual specimen handling and test
setup.

We define a parameterized family of unit cells using a core-wrapping
architecture in which rigid PLA struts are enclosed by continuous TPU
skins to improve interfacial interlocking.  Design variables include
strut diameter and length, tension-element width and thickness, strut
count per unit cell, and tiling pattern.  Specimens are fabricated on a
multi-material FDM system and tested under quasi-static compression and
instrumented drop-weight impact.

We will report how BO-selected designs evolve across iterations, how the
surrogate model changes as data accumulate, and what Pareto-efficient
trade-offs emerge between SEA and peak transmitted force within the
explored design space.  We will also discuss practical lessons for
operating the closed-loop workflow, including print-failure handling,
batch-to-batch variability, and the exploration--exploitation trade-off.

\end{document}
